\documentclass[10pt,conference,letterpaper]{IEEEtran}
\providecommand{\DoNotLoadEpstopdf}{}
\usepackage[T1]{fontenc}
\usepackage{lmodern,amsmath,amssymb,graphicx,booktabs}
\usepackage[american]{circuitikz}
% Shared drawing settings; draw_atlas.py embeds these in standalone exports.
\ctikzset{transistors/scale=0.8, resistors/scale=0.5, capacitors/scale=0.5, sources/scale=0.6}
\ctikzset{logic ports/scale=0.6}
\ctikzset{flipflops/scale=0.6}
\ctikzset{sources/scale=0.8}
\ctikzset{tripoles/mos style=arrows}
\ctikzset{transistors/arrow pos=end}
\ctikzset{amplifiers/scale=0.7}
\ctikzset{blocks/scale=0.7}

\usepackage{pgfplots}
\pgfplotsset{compat=1.17}
\usetikzlibrary{calc,positioning,arrows.meta,shapes.gates.logic.US}
\usepackage[hidelinks]{hyperref}
\hypersetup{pdftitle={Coarse Acquisition and Analog Handoff in a 3.2 GHz Ring PLL},pdfauthor={PLL130-AFC Project}}
\title{Coarse Acquisition and Analog Handoff in a\\3.2\,GHz Ring PLL in 130\,nm CMOS}
\author{\IEEEauthorblockN{PLL130-AFC Project}\IEEEauthorblockA{Design and verification report}}
\begin{document}
\maketitle
\begin{abstract}
A five-bit calibration controller places a two-stage ring oscillator within the capture range of a 100\,MHz-reference charge-pump PLL. Two switches hold both filter nodes during search. A nominal transistor-level transient at 1.2\,V and 25\,$^\circ$C selected code 29 and met the acquisition criteria 4.690\,$\mu$s after calibration enable. The last 500\,ns consumed 4.185\,mW. A second run reached fine lock but exceeded its precharge-error limit. A fixed-code control remained below target throughout the 6.5\,$\mu$s observation. The schematics connect these outcomes to the implemented current bank, charge pump, clamp and controller.
\end{abstract}
\begin{IEEEkeywords}
automatic frequency calibration, charge-pump PLL, ring oscillator, coarse tuning, transistor-level simulation
\end{IEEEkeywords}
\section{Design Objective}
The loop must multiply a 100\,MHz reference by 32. Its control voltage can tune the oscillator only within the band selected by the current bank. If that band is too low, the control voltage can approach the supply before the divided clock reaches the reference frequency. The fixed-code experiment in Section~\ref{sec:result} exhibits this behavior.

The starting 180\,nm project describes two gated counts, a numeric frequency error $E$, and a band update $S\leftarrow S+E$ \cite{student}. Its supplied two-stage VCO files informed the ring topology. Its separate four-stage top-level SPICE file instantiates a conventional PLL. The present 130\,nm controller adapts the two-phase direction decision and bounded search of Jeong \emph{et al.} \cite{jeong}. Their LC oscillator uses a capacitor bank. Here the search controls ring current branches, and the two observation phases come from the reference clock and its inverse. Five static CMOS toggle stages divide the clock. Drain-switched mirrors replace the regulated source-switched pump shown in the student report.
\section{Signal Path and Handoff}
The physical feedback clock closes the analog loop and supplies the AFC measurement input (Fig.~\ref{fig:system}). During coarse search the pump inputs are disabled. Transmission gates connect CTRL and the midpoint of the filter branch to a common half-supply reservoir. Holding both nodes limits the initial voltage difference across the filter resistor at handoff. Once the controller accepts a band, it retains the code, releases the clamp and permits the PFD pulses to reach the pump.
\begin{figure*}[!t]\centering
\begin{circuitikz}[font=\small,blk/.style={draw,align=center,minimum height=9mm,minimum width=18mm},arr/.style={-{Latex[length=2mm]}}]
\node[blk] (pfd) at (1.6,0) {PFD}; \node[blk] (cp) at (4.3,0) {Enable gates\\and pump};
\node[blk] (lf) at (7.4,0) {Loop filter}; \node[blk] (vco) at (10.2,0) {Two-stage\\ring VCO};
\node[blk] (rx) at (13.5,0) {AC receiver}; \node[blk] (div) at (12,-1.55) {CMOS $\div 32$};
\node[blk] (afc) at (7.4,2) {Two-phase AFC\\and code register};
\draw[arr] (-.25,0) node[left]{REF} -- (pfd.west); \draw[arr] (pfd) -- node[above,font=\tiny,align=center,inner sep=1pt]{UP\\DN} (cp);
\draw[arr] (cp) -- (lf); \draw[arr] (lf) -- node[above,font=\scriptsize]{CTRL} (vco);
\draw[arr] (vco) -- node[above,font=\scriptsize]{ring output} (rx);
\draw[arr] (rx.south) |- (div.east); \draw[arr] (div.west) -| (pfd.south);
\draw[arr] (.18,0) node[circ]{} |- (afc.west);
\draw[arr] (afc.east) -| node[pos=.38,above,font=\scriptsize]{five-bit code} (vco.north);
\draw[arr] (afc.south) -- node[right,font=\scriptsize]{hold} (lf.north);
\draw[arr] (6.7,1.55) -- (6.7,1.02) -| (cp.north);
\node[font=\scriptsize,anchor=south] at (5.3,1.04) {fine enable};
\draw (8.5,-1.55) node[circ]{} -- ++(0,-.25) node[below,font=\scriptsize]{FB};
\draw[arr] (8.1,2.95) node[above,font=\scriptsize]{FB} -- (8.1,2.45);
\end{circuitikz}
\caption{Selected PLL. Both FB labels identify the physical divided clock. AFC also supplies PFD reset. Each clamp connects one filter node to the half-supply reservoir.}\label{fig:system}\end{figure*}
\section{Acquisition Experiment}\label{sec:result}
Spectre 18.1.0.077 simulated the selected netlist at the typical corner, 25\,$^\circ$C and 1.2\,V. The bench used ideal reference and bias sources, initial conditions for the ring and filter, a 4\,ps maximum integration step and a 40\,ps saved strobe. The deterministic 6.5\,$\mu$s observation contains 162\,501 points. Values below come from the archived runs \cite{evidence}.

The reader interpolated rising zero crossings of $v_0-v_{180}$. For $N$ edges at times $t_1,\ldots,t_N$ in a 100\,ns window it computed $f_{\mathrm{win}}=(N-1)/(t_N-t_1)$. Windows advanced by 50\,ns. Acquisition required frequency error within 0.1\% of target, unwrapped feedback/reference timing variation within 100\,ps peak to peak over a 500\,ns hold, and output swing of at least 50\,mV peak to peak. Every subsequent complete hold window also had to pass. A separate handoff criterion limited the held CTRL error to 30\,mV relative to 0.6\,V.
\begin{table}[!t]\centering\caption{Frozen counts from the selected calibration}\label{tab:codes}
\begin{tabular}{rrrrl}\toprule
Code & Lead & Lag & Feedback & Decision\\\midrule
16&9&9&7&Slow\\24&25&24&22&Slow\\28&81&81&79&Slow\\30&67&67&69&Fast\\29&130&128&129&Near\\\bottomrule
\end{tabular}\end{table}
The first three decisions increased the coarse code. At code 30 the feedback count exceeded the lead count, so the last step selected 29. DONE rose at 3.850960\,$\mu$s and FineEnable at 3.870949\,$\mu$s. The final mean CTRL was 0.64750\,V and differential swing was 0.22705\,V peak to peak.

\begin{figure}[!t]\centering\begin{tikzpicture}
\begin{axis}[width=\columnwidth,height=4.6cm,xmin=0,xmax=6.5,ymin=2.1,ymax=3.6,xlabel={Time from simulation start ($\mu$s)},ylabel={$f_{\mathrm{win}}$ (GHz)},tick label style={font=\footnotesize},label style={font=\footnotesize},grid=major,legend style={font=\scriptsize,at={(.97,.04)},anchor=south east}]
\addplot[blue!65!black,thick] table[x=time_us,y=on_GHz]{acquisition.dat};\addlegendentry{AFC enabled}
\addplot[black,dashed,thick] table[x=time_us,y=off_GHz]{acquisition.dat};\addlegendentry{Fixed code 16}
\addplot[gray,densely dotted,forget plot] coordinates {(0,3.2) (6.5,3.2)};
\end{axis}\end{tikzpicture}
\caption{Output frequency at the centers of all 129 complete windows in each archived trace. The dotted line marks 3.2\,GHz.}\label{fig:freq}\end{figure}
\begin{table}[!t]\centering\caption{Nominal closed-loop outcomes}
\begin{tabular}{lrr}\toprule
Quantity & AFC enabled & Fixed code 16\\\midrule
Tail frequency (GHz)&3.199994&2.469028\\
Tail mean CTRL (V)&0.64750&1.18790\\
Tail supply power (mW)&4.18516&3.38261\\
Absolute acquisition ($\mu$s)&4.98955&none in 6.5\\
Enable to acquisition ($\mu$s)&4.68953&---\\\bottomrule
\end{tabular}\end{table}
With calibration disabled, CTRL approached the positive rail while the output remained below target. This control removed the AFC loading and clamp circuitry. Its tail power describes an uncaptured operating point. In the enabled circuit, supply energy from calibration enable to lock was 24.388\,nJ.

The weaker-clamp run reached fine lock but exceeded the 30\,mV precharge limit with 35.547\,mV maximum error. Increasing the switch widths and hold capacitance reduced the error to 19.313\,mV. This comparison measures their combined effect. A separate unloaded VCO run at code 26 with 60 sidebands gave $-73.478$ and $-95.231$\,dBc/Hz at 1 and 10\,MHz offsets. Integrated closed-loop phase noise, RMS jitter, reference spurs and post-layout acquisition remain unverified.
\section{Using the Sources}
The canonical AFC-on bench, fixed-code control, measurement ledgers and window series are available in the repository \cite{evidence}. The Spectre decks define the electrical experiment. OCEAN scripts export noise data from stored results. Project SKILL includes a one-line Library Manager launcher. A local installation supplies the process models. The source archive distinguishes the original simulation input hashes from portable copy hashes.
% Selected R1 implementation. Parent preamble: graphicx, circuitikz[american].
% All figures are vector schematics. Like-named labeled terminals are joined.
\section{Transistor Schematics of the Selected Loop}
\label{sec:schematicatlas}
The following schematics follow the instantiated R1 strong-clamp hierarchy.
The source identifiers in Table~\ref{tab:apmap} can be searched verbatim in
the accompanying netlists. A repeated net label denotes electrical
continuity within its stated module. An NMOS bulk connects to its local
negative rail and a PMOS bulk to its local positive rail. The captions identify RF devices. Each of these symbols represents a four-terminal RF device wrapper at the design boundary.
Its finger count belongs to one instance. All component values refer to the
selected pre-layout bench.

\begin{table*}[!t]
\centering\caption{Schematic-to-netlist map for the selected loop.}
\label{tab:apmap}
\small
\resizebox{.99\textwidth}{!}{\begin{tabular}{lll}
\hline
Selected instance & Netlist module or devices & Function \\
\hline
\texttt{XREF}, \texttt{XOSC} & \texttt{current\_ref, two\_stage\_ring\_vco, two\_stage\_ring} & Bias and two quadrature cells \\
\texttt{XRX} & \texttt{pll130\_receiver\_ac} & Differential input to logic clock \\
\texttt{XD} & \texttt{pll130\_div32\_cmos, pll130\_div2, pll130\_dff, pll130\_latch} & Five toggle stages \\
\texttt{XPFD} & \texttt{pll130\_pfd, pll130\_pfd\_tspc} & Reference and feedback phase comparison \\
\texttt{XCPUPEN}, \texttt{XCPDNEN}, \texttt{XCP} & \texttt{pll130\_and2, pll130\_cp\_drain, pll130\_cpd\_inv} & Pulse gating and charge transfer \\
Top-level passives and switches & \texttt{CSHUNT, RSERIES, CSERIES, CHOLD, MHOLD*, PHOLD*} & Integration and precharge \\
\hline
\end{tabular}}
\end{table*}


\tikzset{apbox/.style={draw,line width=.85pt,align=center,minimum height=6mm,inner xsep=4pt,font=\PLLfigurefont\footnotesize},
apnote/.style={align=left,font=\PLLfigurefont\footnotesize},
aphead/.style={font=\PLLfigurefont\small,anchor=west}}
\newcommand{\APgleft}[2]{\draw (#1.G) -- ++(-.24,0) node[left]{#2};}
\newcommand{\APgright}[2]{\draw (#1.G) -- ++(.24,0) node[right]{#2};}
\newcommand{\APtop}[2]{\draw (#1) -- ++(0,.27) node[above]{#2};}
\newcommand{\APbot}[2]{\draw (#1) -- ++(0,-.27) node[below]{#2};}

\subsection{Oscillator and band selection}
The oscillator closes two differential stages with one polarity swap.
Stage~0 takes $(o_0,o_{180})$ and drives $(o_{90},o_{270})$.
Stage~1 takes $(o_{270},o_{90})$ and drives $(o_0,o_{180})$.
In each stage the input pair and the regenerative pair have separate tail
nodes (Fig.~\ref{fig:apcore}). The control branches act on the input-pair tail
only (Fig.~\ref{fig:apbank}). The selected sizing is given by the actual
instance parameters. In particular the RF input devices have twelve fingers
of width $3\,\mu$m each.

\begin{figure*}[!t]
\centering\resizebox{.92\textwidth}{!}{%
\begin{circuitikz}[font=\PLLfigurefont\footnotesize]
\node[aphead] at (-1.4,5.35) {(a) Main differential pair};
\node[nmos] (am1) at (0,2.75) {};
\node[nmos,xscale=-1] (am2) at (2.5,2.75) {};
\draw (0,4.85) node[above]{$V_{DD}$} to[R,l=$620\,\Omega$] (0,3.55) node[circ]{} -- (am1.D);
\draw (2.5,4.85) node[above]{$V_{DD}$} to[R,l=$620\,\Omega$] (2.5,3.55) node[circ]{} -- (am2.D);
\draw (0,3.55) -- ++(-.45,0) node[left]{$o_n$};
\draw (2.5,3.55) -- ++(.45,0) node[right]{$o_p$};
\APgleft{am1}{$i_p$}\APgright{am2}{$i_n$}
\draw (am1.S) -- (0,1.95) -- (2.5,1.95) -- (am2.S);
\node[nmos] (amt) at (1.25,.95) {};\APgleft{amt}{$b_n$}
\draw (amt.D) -- (1.25,1.95) node[circ]{} node[above]{$c_{s1}$};
\draw (amt.S) node[ground]{};
\draw (2.5,1.95) -- ++(.65,0) node[right]{to fine/bank};
\begin{scope}[xshift=6.6cm]
\node[aphead] at (-1.4,5.35) {(b) Regenerative pair};
\node[nmos] (ar1) at (0,2.75) {};
\node[nmos,xscale=-1] (ar2) at (2.25,2.75) {};
\APtop{ar1.D}{$o_n$}\APtop{ar2.D}{$o_p$}
\APgleft{ar1}{$o_p$}\APgright{ar2}{$o_n$}
\draw (ar1.S) -- (0,1.95) -- (2.25,1.95) -- (ar2.S);
\node[nmos] (art) at (1.125,.95) {};\APgleft{art}{$b_n$}
\draw (art.D) -- (1.125,1.95) node[circ]{} node[above]{$c_{s2}$};
\draw (art.S) node[ground]{};
\end{scope}
\begin{scope}[xshift=12cm]
\node[aphead] at (-1.15,5.35) {(c) Bias reference};
\node[nmos] (anref) at (0,1.25) {};
\node[nmos] (anmir) at (2.5,1.25) {};
\node[pmos] (apref) at (2.5,3.3) {};
\draw (0,4.85) node[above]{$V_{DD}$} to[I,l=$I_{\rm REF}$] (0,2.35) node[circ]{} -- (anref.D);
\draw (anref.G) -- ++(-.18,0) |- (0,2.35) node[above right]{$b_n$};
\draw (anref.S) node[ground]{};\draw (anmir.S) node[ground]{};
\APgleft{anmir}{$b_n$}
\draw (apref.S) -- ++(0,.36) node[above]{$V_{DD}$};
\draw (apref.D) -- (2.5,2.3) node[circ]{} -- (anmir.D);
\draw (apref.G) -- ++(-.2,0) |- (2.5,2.3) node[right]{$b_p$};
\draw (4.15,2.3) node[above]{$b_n$} to[C,l=$1\,$pF] (4.15,.8) node[ground]{};
\end{scope}
\end{circuitikz}}
\caption{One ring cell and the common bias source. The six devices in (a,b) use RF wrappers. Input, regenerative, main-tail and regenerative-tail finger counts are 12, 8, 12 and 16, respectively. Each finger has $W=3\,\mu$m. Pair devices have $L=0.13\,\mu$m and tail devices $L=0.39\,\mu$m. The bias source supplies $13.5594\,\mu$A. The two bias NMOS devices have $W/L=8/3.12$ and the PMOS has $W/L=4.06783/0.39$, in $\mu$m.}
\label{fig:apcore}
\end{figure*}

\begin{figure*}[!t]
\centering\resizebox{.91\textwidth}{!}{%
\begin{circuitikz}[font=\PLLfigurefont\footnotesize]
\node[aphead] at (-1.2,4.35) {(a) Fine tuning};
\node[nmos] (af) at (0,1.7) {};
\APgleft{af}{CTRL}\APtop{af.D}{$c_{s1}$}\draw (af.S) node[ground]{};
\begin{scope}[xshift=3.45cm]
\node[aphead] at (-1,4.35) {(b) Binary and correction branches};
\node[nmos] (ab1) at (0,2.45) {};\node[nmos,xscale=-1] (ab2) at (2.2,2.45) {};
\APtop{ab1.D}{$g_j$}\APtop{ab2.D}{$g_j$}
\APgleft{ab1}{$e_j$}\draw (ab2.G) -- ++(.18,0) node[above]{$\bar e_j$};
\APbot{ab1.S}{$b_n$}\draw (ab2.S) node[ground]{};
\node[nmos,xscale=-1] (ab3) at (4.4,2.45) {};\APgright{ab3}{$g_j$}\APtop{ab3.D}{$c_{s1}$}\draw (ab3.S) node[ground]{};
\node[apnote,anchor=west] at (-.85,.45) {$j=0\ldots4$: $e_j=s_j$, $m=1,2,4,8,16$\\$j=5,6$: $e_j=s_{24+},s_{16+}$};
\end{scope}
\begin{scope}[xshift=11.2cm]
\node[aphead] at (-1.3,4.35) {(c) Supplemental fine branches};
\node[nmos] (as1) at (0,2.45) {};\node[pmos,xscale=-1,yscale=-1] (as2) at (1.8,2.45) {};
\APgleft{as1}{$e_k$}
\draw (as1.D) -- (0,3.3) -- (1.8,3.3) -- (as2.D);\draw (.9,3.3) node[circ]{} -- ++(0,.2) node[above]{$g_k$};
\draw (as1.S) -- (0,1.6) -- (1.8,1.6) -- (as2.S);\draw (.9,1.6) node[circ]{} -- ++(0,-.15) node[below]{CTRL};
\node[nmos] (as3) at (3.9,2.45) {};\draw (as2.G) -- node[above]{$\bar e_k$} (as3.G);\APtop{as3.D}{$g_k$}\draw (as3.S) node[ground]{};
\node[nmos] (as4) at (6.1,2.45) {};\draw (as4.G) -- ++(-.2,0) node[below]{$g_k$};\APtop{as4.D}{$c_{s1}$}\draw (as4.S) node[ground]{};
\node[apnote,anchor=west] at (-1,.45) {$e_k=s_{8+},s_{16+},s_{24+}$; supplemental $m=1,2,3$};
\end{scope}
\end{circuitikz}}
\caption{All three tuning paths in one cell. The fine device has $W/L=0.8/5.2\,\mu$m. Binary sinks have unit $W/L=1.1/0.39$. The two correction widths are 0.47143 and 0.23571\,$\mu$m for $24+$ and $16+$, respectively. For (c), the selected device gate receives CTRL through a transmission gate and the disabled gate is grounded. Supplemental devices have $W/L=0.2/5.2$. The decoder forms $s_{8+}=s_3\lor s_4$, $s_{16+}=s_4$ through two inverters and $s_{24+}=s_3\land s_4$. The \texttt{s\_to\_d} buffers in Fig.~\ref{fig:apdecode} derive both polarities. The labels $g_j$ and $g_k$ refer to separate branch families.}
\label{fig:apbank}
\end{figure*}

\begin{figure*}[!t]
\centering\resizebox{.67\textwidth}{!}{%
\begin{circuitikz}[font=\PLLfigurefont\footnotesize]
\node[aphead] at (-.65,3.1) {(a) Complementary drive\ \texttt{s\_to\_d}};
\node[american not port] (sd1) at (2,2) {};
\node[american not port] (sd2) at (4.65,2) {};
\node[apbox] (dtg) at (2,.35) {TG (always on)};
\node[american not port] (sd3) at (4.65,.35) {};
\draw (-.5,2) node[left]{$s$} -- (.3,2) node[circ]{} -- (sd1.in);
\draw (.3,2) |- (dtg.west);\draw (sd1.out) -- (sd2.in);
\draw (dtg.east) -- (sd3.in);
\draw (sd2.out) -- ++(.3,0) node[right]{$s_p$};
\draw (sd3.out) -- ++(.3,0) node[right]{$s_n$};
\begin{scope}[xshift=8.0cm]
\node[aphead] at (-.6,3.1) {(b) Decode and drive};
\node[american or port] (bor) at (1.5,2) {};
\node[american and port] (band) at (1.5,-.25) {};
\draw (bor.in 1) -- ++(-.25,0) node[left]{$s_3$};\draw (bor.in 2) -- ++(-.25,0) node[left]{$s_4$};
\draw (band.in 1) -- ++(-.25,0) node[left]{$s_3$};\draw (band.in 2) -- ++(-.25,0) node[left]{$s_4$};
\draw (bor.out) -- ++(.55,0) node[right]{$s_{8+}$};
\draw (band.out) -- ++(.55,0) node[right]{$s_{24+}$};
\node[apbox] (b16) at (1.55,.875) {\texttt{bf\_2inv}};
\draw (b16.west) -- ++(-.3,0) node[left]{$s_4$};
\draw (b16.east) -- ++(.6,0) node[right]{$s_{16+}$};
\end{scope}
\end{circuitikz}}
\caption{VCO bank logic. Each cell contains eight complementary drivers, one per binary bit and one per decoded threshold. The always-on TG connects the input to a single inverter with its NMOS gate at VDD and PMOS gate at VSS. The other branch uses two inverters. The \texttt{bf\_2inv} buffer is also a pair of inverters in series. OR2 uses NOR2 followed by an inverter, and AND2 uses NAND2 followed by an inverter. All VCO logic inverter devices have $W_n/W_p=0.5/1.8\,\mu$m and $L=0.13\,\mu$m.}
\label{fig:apdecode}
\end{figure*}
\subsection{Clock conversion and frequency division}
The receiver couples the positive differential-pair output to a CMOS clock
restorer (Fig.~\ref{fig:aprx}). Feedback biases the first inverter near its
switching threshold. Five toggle stages then divide the recovered clock by
32 (Fig.~\ref{fig:apdiv}).

\begin{figure*}[!t]
\centering\resizebox{.80\textwidth}{!}{%
\begin{circuitikz}[font=\PLLfigurefont\footnotesize]
\node[nmos] (ai1) at (0,2.0) {};\node[nmos,xscale=-1] (ai2) at (3.1,2.0) {};
\draw (0,4.5) node[above]{$V_{DD,\rm dig}$} to[R,l=$3\,$k$\Omega$] (0,3.05) node[circ]{} -- (ai1.D);
\draw (3.1,4.5) node[above]{$V_{DD,\rm dig}$} to[R,l=$3\,$k$\Omega$] (3.1,3.05) node[circ]{} -- (ai2.D);
\draw (0,3.05) -- ++(-.45,0) node[left]{rxn};
\APgleft{ai1}{$o_0$}\APgright{ai2}{$o_{180}$}
\draw (ai1.S) -- (0,1.15) -- (3.1,1.15) -- (ai2.S);
\node[nmos] (ait) at (1.55,.05) {};\APgleft{ait}{$b_n$}
\draw (ait.D) -- (1.55,1.15) node[circ]{};\draw (ait.S) node[ground]{};
\node[american not port] (axi) at (8.0,3.05) {};
\node[american not port] (axi2) at (10.5,3.05) {};
\draw (3.1,3.05) -- (4.3,3.05) node[above]{rxp} to[C,l=$40\,$fF] (6.2,3.05) node[circ]{} -- (axi.in);
\draw (axi.out) -- (9.1,3.05) node[circ]{} -- (axi2.in);
\draw (9.1,3.05) -- (9.1,4.35) to[R,l=$1\,$M$\Omega$] (6.2,4.35) -- (6.2,3.05);
\draw (axi2.out) -- ++(.45,0) node[right]{clk};
\node[apnote,anchor=north] at (8.7,1.95) {Stage 1 $W_n/W_p=0.8/3.2\,\mu$m\\Stage 2 $W_n/W_p=1.6/4.8\,\mu$m};
\end{circuitikz}}
\caption{Selected AC-coupled receiver. Each differential input uses four RF fingers of $1.5\,\mu$m at $L=0.13\,\mu$m. The tail uses 12 fingers of $2.5\,\mu$m at $L=0.39\,\mu$m. The bench sets the coupling capacitance to 40 fF.}
\label{fig:aprx}
\end{figure*}

\begin{figure*}[!t]
\centering\resizebox{.94\textwidth}{!}{%
\begin{circuitikz}[font=\PLLfigurefont\footnotesize]
\node[aphead] at (-.5,4.8) {(a) Transmission-gate latch};
\node[apbox] (dIn) at (1.8,3.6) {TG$_{\rm open}$};
\node[american not port] (dX) at (5.5,3.6) {};
\node[american not port] (dY) at (8.35,3.6) {};
\draw (-.1,3.6) node[left]{D} -- (dIn.west);
\draw (dIn.east) -- (3.5,3.6) node[circ]{} node[above]{storage} -- (dX.in);
\draw (dX.out) -- (6.65,3.6) node[circ]{} -- (dY.in);
\draw (6.65,3.6) -- ++(0,.5) node[above]{$\bar Q$};
\draw (dY.out) -- (9.65,3.6) node[circ]{} -- ++(.3,0) node[right]{$Q$};
\node[apbox] (dFb) at (6.4,1.55) {TG$_{\rm hold}$};
\draw (9.65,3.6) |- (dFb.east);\draw (dFb.west) -| (3.5,3.6);
\node[nmos] (dR) at (2.25,.7) {};\APgleft{dR}{rst}\draw (dR.S) node[ground]{};
\draw (dR.D) |- (3.5,2.25) node[circ]{};
\begin{scope}[xshift=12cm]
\node[aphead] at (-.8,4.8) {(b) Each TG};
\node[nmos] (dTn) at (0,2.7) {};\node[pmos,xscale=-1,yscale=-1] (dTp) at (1.65,2.7) {};
\APgleft{dTn}{$E$}\APgright{dTp}{$\bar E$}
\draw (dTn.D) -- (0,3.6) -- (1.65,3.6) -- (dTp.D);\draw (.825,3.6) node[circ]{} -- ++(0,.15) node[above]{A};
\draw (dTn.S) -- (0,1.8) -- (1.65,1.8) -- (dTp.S);\draw (.825,1.8) node[circ]{} -- ++(0,-.15) node[below]{B};
\end{scope}
\node[aphead] at (-.5,-1.25) {(c) Positive-edge toggle stage};
\node[apbox] (dM) at (3.55,-2.3) {Master latch\\open $=L$, hold $=H$};
\node[apbox] (dS) at (8.4,-2.3) {Slave latch\\open $=H$, hold $=L$};
\draw[->] (0,-2.3) node[left]{$\bar Q$} -- (dM.west);
\draw[->] (dM.east) -- node[above]{$Q_M$} (dS.west);
\draw[->] (dS.east) -- ++(.75,0) node[right]{$Q,\bar Q$};
\node[apnote,anchor=west] at (12.2,-2.3) {$L=\overline{\rm clk}\land\overline{\rm rst}$\\$H={\rm clk}\land\overline{\rm rst}$};
\end{circuitikz}}
\caption{Physical divider hierarchy. Five copies of (c) cascade through their $Q$ outputs, giving $q_1,q_2,q_3,q_4,\mathrm{fb}$. A $5$-fF capacitor loads fb. Each latch has two TGs, two inverters and one reset NMOS. TGs have $W_n/W_p=1.2/3.6\,\mu$m, latch inverters $0.5/1.5\,\mu$m and the reset NMOS $W=2\,\mu$m. Inverters outside the latch generate complementary TG controls.}
\label{fig:apdiv}
\end{figure*}

\subsection{Phase detector, charge pump, and precharge}
The two six-device PFD slices receive REF and divider feedback
(Fig.~\ref{fig:appfd}). After the controller enables fine tuning, two AND
gates pass UP and DN to the charge pump. The pump switches the drains of
continuously biased mirror devices (Fig.~\ref{fig:apcp}). During calibration
a transmission gate at each filter node connects that node to the half-supply
reservoir (Fig.~\ref{fig:apclamp}).

\begin{figure*}[!t]
\centering\resizebox{.84\textwidth}{!}{%
\begin{circuitikz}[font=\PLLfigurefont\footnotesize]
\node[aphead] at (-1.4,6.65) {(a) One PFD slice};
\node[pmos] (pA) at (0,4.95) {};\node[pmos] (pB) at (0,3.45) {};\node[nmos] (pC) at (0,1.55) {};
\APtop{pA.S}{$V_{DD}$}\APgleft{pA}{reset}\APgleft{pB}{edge}\APgleft{pC}{reset}
\draw (pA.D) -- (pB.S);\draw (pB.D) -- (0,2.45) node[circ]{} -- (pC.D);\draw (0,2.45) -- ++(.6,0) node[right]{$n_2$};\draw (pC.S) node[ground]{};
\node[pmos] (pD) at (3.55,4.95) {};\node[nmos] (pE) at (3.55,3.1) {};\node[nmos] (pF) at (3.55,1.55) {};
\APtop{pD.S}{$V_{DD}$}\APgleft{pD}{$n_2$}\APgleft{pE}{edge}\APgleft{pF}{$n_2$}
\draw (pD.D) -- (3.55,4.0) node[circ]{} -- (pE.D);\draw (3.55,4.0) -- ++(.45,0) node[right]{$y$};\draw (pE.S) -- (pF.D);\draw (pF.S) node[ground]{};
\begin{scope}[xshift=7.25cm]
\node[aphead] at (-.9,6.65) {(b) Pair and reset logic};
\node[apbox] (pu) at (1.35,4.95) {PFD slice};
\node[apbox] (pd) at (1.35,3.4) {PFD slice};
\draw[->] (-.25,4.95) node[left]{REF} -- (pu.west);
\draw[->] (-.25,3.4) node[left]{FB} -- (pd.west);
\node[american not port] (xu) at (4.8,4.95) {};
\node[american not port] (xd) at (4.8,3.4) {};
\draw (pu.east) -- node[above]{$u_n$} (xu.in);\draw (xu.out) -- ++(.35,0) node[right]{UP};
\draw (pd.east) -- node[above]{$d_n$} (xd.in);\draw (xd.out) -- ++(.35,0) node[right]{DN};
\node[apnote,anchor=west] at (-.75,1.6) {$r_0=\mathrm{NOR}(u_n,d_n)$\\$r_2=\mathrm{INV}(\mathrm{INV}(r_0))$\\$\mathrm{reset}=\mathrm{pfd\_reset}\lor r_2$};
\end{scope}
\end{circuitikz}}
\caption{Six-MOS single-phase dynamic slice (left) and PFD assembly (right). The two PMOS devices in the input column are stacked. The reset NMOS discharges $n_2$. In the output column $n_2$ drives the PMOS and bottom NMOS, while the edge input drives the middle NMOS. The two explicit inverters preserve the delays of the reset path.}
\label{fig:appfd}
\end{figure*}

\begin{figure*}[!t]
\centering\resizebox{.81\textwidth}{!}{%
\begin{circuitikz}[font=\PLLfigurefont\footnotesize]
\node[aphead] at (-1.2,6.65) {(a) Pump bias};
\node[nmos] (cN) at (0,1.25) {};\node[nmos] (cN2) at (3.1,1.25) {};
\node[pmos] (cP) at (3.1,4.3) {};
\draw (0,5.05) node[above]{$V_{DD,\rm CP}$} to[I,l=$20\,\mu$A] (0,2.42) node[circ]{} -- (cN.D);
\draw (cN.G) -- ++(-.25,0) |- (0,2.42) node[right]{$b_{nC}$};
\draw (cN.S) node[ground]{};\draw (cN2.S) node[ground]{};
\APgleft{cN2}{$b_{nC}$}
\APtop{cP.S}{$V_{DD,\rm CP}$}
\draw (cP.D) -- (3.1,2.92) node[circ]{} -- (cN2.D);
\draw (cP.G) -- ++(-.20,0) |- (3.1,2.92) node[right]{$b_{pC}$};
\begin{scope}[xshift=8cm]
\node[aphead] at (-1.8,6.65) {(b) Output branch};
\node[pmos] (cM1) at (0,4.75) {};
\node[pmos] (cM2) at (0,3.35) {};
\node[nmos] (cM3) at (0,1.65) {};
\node[nmos] (cM4) at (0,.25) {};
\APtop{cM1.S}{$V_{DD,\rm CP}$}
\draw (cM1.D)--(cM2.S);\draw (cM2.D)--(0,2.5) node[circ]{}--(cM3.D);
\draw (0,2.5)--++(.85,0) node[right]{CTRL};\draw (cM3.S)--(cM4.D);\draw (cM4.S) node[ground]{};
\APgleft{cM1}{$b_{pC}$}\APgleft{cM2}{$\overline{u_g}$}\APgleft{cM3}{$d_g$}\APgleft{cM4}{$b_{nC}$}
\node[apnote,anchor=west] at (1.3,4.5) {$u_g=\mathrm{UP}\land\mathrm{FineEnable}$\\$d_g=\mathrm{DN}\land\mathrm{FineEnable}$};
\end{scope}
\end{circuitikz}}
\caption{Drain-switched charge pump. Bias mirrors have $W_n/L=4/1.04$ and $W_p/L=12/1.04$ in $\mu$m. Output switches have $W_n/L=2/0.13$ and $W_p/L=6/0.13$. A CMOS inverter generates $\overline{u_g}$. The separate oscillator bias of Fig.~\ref{fig:apcore} has its own bias nodes and reference current.}
\label{fig:apcp}
\end{figure*}

\begin{figure*}[!t]
\centering\resizebox{.90\textwidth}{!}{%
\begin{circuitikz}[font=\PLLfigurefont\footnotesize]
\node[aphead] at (-1.2,5.7) {(a) Half-supply reservoir};
\draw (0,4.75) node[above]{$V_{DD,\rm AFC}$} to[R,l=$20\,$k$\Omega$] (0,2.6) node[circ]{} to[R,l=$20\,$k$\Omega$] (0,.65) node[ground]{};
\draw (0,2.6) -- (2,2.6) node[circ]{} node[above]{vhold} to[C,l=$10\,$pF] (2,.65) node[ground]{};
\begin{scope}[xshift=5.4cm]
\node[aphead] at (-.6,5.7) {(b) Filter and clamps};
\draw (0,3) node[circ]{} node[above left]{CTRL} to[R,l=$49.5015\,$k$\Omega$] (3.7,3) node[circ]{} node[above right]{filter\_mid};
\draw (0,3) to[C,l_=$0.338437\,$pF] (0,.95) node[ground]{};
\draw (3.7,3) to[C,l=$6.43031\,$pF] (3.7,.95) node[ground]{};
\node[apbox] (ch1) at (.0,4.05) {TG$_H$};\node[apbox] (ch2) at (3.7,4.05) {TG$_H$};
\draw (0,3) -- (ch1.south);\draw (3.7,3) -- (ch2.south);
\draw (ch1.north) -- (0,5.0) -- (3.7,5.0) -- (ch2.north);
\draw (1.85,5.0) node[circ]{} node[above]{vhold};
\end{scope}
\begin{scope}[xshift=12.9cm]
\node[aphead] at (-1.4,5.7) {(c) Each clamp TG};
\node[nmos] (hN) at (0,2.9) {};\node[pmos,xscale=-1,yscale=-1] (hP) at (1.75,2.9) {};
\APgleft{hN}{$H$}\APgright{hP}{$\bar H$}
\draw (hN.D) -- (0,3.9) -- (1.75,3.9) -- (hP.D);\draw (.875,3.9) node[circ]{} node[above]{filter node};
\draw (hN.S) -- (0,1.9) -- (1.75,1.9) -- (hP.S);\draw (.875,1.9) node[circ]{} node[below]{vhold};
\end{scope}
\end{circuitikz}}
\caption{Both filter nodes precharge from the same passive reference during coarse calibration. The complementary control pair is $H=\mathrm{hold\_all}$ and $\bar H=\mathrm{INV}(H)$. Each physical transmission gate uses a $16/0.13\,\mu$m NMOS and a $32/0.13\,\mu$m PMOS. These ratios identify the strong-clamp revision.}
\label{fig:apclamp}
\end{figure*}

% Inline source for IEEEtran. Requires graphicx, amsmath and american circuitikz.
\section{Transistor Implementation of the AFC}
\label{sec:r1afc}
The R1 AFC compares edges of the 100-MHz reference and the physical
divide-by-32 feedback. A higher five-bit code increases the ring frequency.
The controller starts at code 16 and can test four further candidates
using steps of 8, 4, 2, and 1. Each candidate reaches the detector before
the controller can hand it to the fine loop.

\begin{figure*}[t]
\centering\resizebox{0.98\textwidth}{!}{%
\begin{circuitikz}[american,x=1cm,y=1cm,font=\small,
 blk/.style={draw,rounded corners=1.5pt,align=center,minimum width=2.5cm,
 minimum height=.76cm,inner sep=4pt}, ar/.style={->,>=latex}]
  \node[blk] (ep) at (1.7,2.2) {First FB edge\\sample REF level};
  \node[blk] (sel) at (5.1,2.2) {Phase selector\\REF or inverted REF};
  \node[blk] (count) at (8.8,2.2) {Lead / lag counters\\8-bit Gray, init.\ 1 / 0};
  \node[blk] (fbg) at (1.7,.55) {FB gate\\enable latched low};
  \node[blk] (div) at (5.1,.55) {FB counter\\8-bit Gray, init.\ 0};
  \node[blk] (sync) at (12.3,1.38) {Two DFFs per bit\\then Gray decode};
  \node[blk] (cmp) at (12.3,-.7) {Compare counts\\qualify direction};
  \node[blk] (snap) at (8.8,-.7) {Decision snapshot\\$3\times8$ bits + direction};
  \node[blk] (fsm) at (5.1,-.7) {Trial controller\\state and timers};
  \node[blk] (code) at (1.7,-.95) {Five-bit band register\\steps $8,4,2,1$};
  \draw[ar] (-.95,2.2) node[left]{REF} -- (ep.west);
  \draw[ar] (-.95,.55) node[left]{FB} -- (fbg.west);
  \draw[ar] (ep.east) -- (sel.west); \draw[ar] (sel.east) -- (count.west);
  \draw[ar] (fbg.east) -- (div.west);
  \draw[ar] (3.26,.55) node[circ]{} -- (3.26,1.32) -| (ep.south);
  \draw[ar] (count.east) -- ++(.20,0) |- ([yshift=2.5mm]sync.west);
  \draw[ar] (div.east) -- (10.35,.55) |- ([yshift=-2.5mm]sync.west);
  \draw[ar] (sync.south) -- (cmp.north); \draw[ar] (cmp.west) -- (snap.east);
  \draw[ar] (snap.west) -- (fsm.east); \draw[ar] (fsm.west) -- ++(-.27,0) |- (code.east);
  \draw[ar] (code.west) -- ++(-.85,0) node[left]{$S[4{:}0]$};
  \draw[ar] (fsm.south) -- ++(0,-.54) node[below]{DONE, FAIL, HOLD, FineEnable, PFD reset};
\end{circuitikz}}
\caption{R1 frequency detector and search controller. The phase selector also receives REF and its complement. The gated FB clocks both the FB counter and the first-edge sample. The synchronized counts share the controller clock. Trial control supplies the counter enables and resets.}
\label{fig:r1afc-datapath}
\end{figure*}

The first admitted FB edge samples the reference level and assigns
the two reference phases to lead and lag counters (Fig.~\ref{fig:r1afc-datapath}).
The lead counter starts with one credit. A chain of 16 inverters guards
reference-counter reset release after that first edge. The controller
samples each Gray word through two flip-flops per bit, then decodes it.
With lead, lag, and divided-feedback counts $C_L,C_G,C_D$, respectively,
it qualifies a slow decision when $C_G-C_D\geq 2$ and a fast decision
when $C_D-C_L\geq 2$. The unsigned subtractors check the carry first.
Two consecutive observations in the same direction qualify a request.
Reaching lag count 128 without a qualified direction ends a trial as near.

\begin{figure}[t]
\centering\resizebox{.93\columnwidth}{!}{%
\begin{circuitikz}[american,x=1cm,y=1cm,font=\small,
 s/.style={draw,rounded corners=1.5pt,minimum width=1.72cm,
 minimum height=.60cm,inner sep=3pt,align=center}, a/.style={->,>=latex}]
 \node[s] (set) at (1,1.5) {SETTLE}; \node[s] (arm) at (4.8,1.5) {ARM};
 \node[s] (meas) at (4.8,0) {MEASURE}; \node[s] (dr) at (1,0) {DRAIN};
 \node[s] (ok) at (1,-1.8) {DONE}; \node[s] (fail) at (4.8,-1.8) {FAIL};
 \draw[a] (1,2.32) node[above]{START} -- (set.north);
 \draw[a] (set.east) -- node[above]{4 REF} (arm.west);
 \draw[a] (arm.south) -- node[right,font=\scriptsize]{epoch valid} (meas.north);
 \draw[a] (meas.west) -- node[above]{decision} (dr.east);
 \draw[a] (dr.west) -- ++(-.45,0) |- ([yshift=-1mm]set.west);
 \node[rotate=90,anchor=south,font=\scriptsize] at (-.33,.75) {next code};
 \draw[a] (dr.south) -- node[left,font=\scriptsize]{accept} (ok.north);
 \draw[a] (meas.south) -- node[right,font=\scriptsize]{watchdog / conflict} (fail.north);
 \draw[a] ([yshift=1mm]arm.east) -- ++(1.62,0) |- (fail.east);
 \node[rotate=90,anchor=south,font=\scriptsize] at (7.15,.5) {ARM timeout};
 \draw[a] ([xshift=4mm]dr.south) -- ++(0,-.91) -| ([xshift=-5mm]fail.north);
 \node[anchor=south,font=\scriptsize] at (2.95,-1.2) {endpoint failure};
\end{circuitikz}}
\caption{R1 state transitions. DRAIN holds the captured counts for two controller cycles. A terminal decision accepts a near result or the last tested candidate. An outward direction at code 1 or 31 on a terminal decision sets FAIL.}
\label{fig:r1afc-fsm}
\end{figure}

An empty ARM window times out after 32 reference clocks. MEASURE has
a 192-clock watchdog. The circuit stores the three decision words and
direction before ending the measurement. On success, two further
reference clocks separate DONE from FineEnable. HOLD is the complement
of FineEnable. The PFD reset is asserted when reset is high or FineEnable
is low.

\subsection{Static CMOS Cells}
The generator expands the gates into one flat Spectre subcircuit.
Figure~\ref{fig:r1afc-cells} shows the transistor connections of the
gate constructors. All devices use $L=0.13~\mu\mathrm{m}$. NMOS bulk
connects to VSS and PMOS bulk connects to VDD. Repeated
input labels in a panel denote the same electrical net.

\begin{figure*}[t]
\centering\resizebox{.97\textwidth}{!}{%
\begin{circuitikz}[american,x=1cm,y=1cm,font=\small]
\begin{scope}
 \node[pmos] (pi) at (0,2.1) {}; \node[nmos] (ni) at (0,.55) {};
 \draw (pi.S) -- (0,2.9) node[above]{VDD};
 \draw (ni.S) -- (0,-.2) node[below]{VSS};
 \draw (pi.D) -- (ni.D); \draw (0,1.325) node[circ]{} -- ++(.65,0) node[right]{$Y$};
 \draw (pi.G) -- (-.84,2.1) -- (-.84,.55) -- (ni.G);
 \draw (-.84,1.325) node[circ]{} -- ++(-.35,0) node[left]{$A$};
 \node at (0,-1.64) {(a) INV};
\end{scope}
\begin{scope}[shift={(3.55,0)}]
 \node[pmos] (pa) at (-.55,2.1) {}; \node[pmos,xscale=-1] (pb) at (.55,2.1) {};
 \node[nmos] (na) at (0,.75) {}; \node[nmos] (nb) at (0,-.28) {};
 \draw (pa.S) -- (-.55,2.92) -- (.55,2.92) -- (pb.S);
 \node[above] at (0,2.92) {VDD};
 \draw (pa.D) -- (-.55,1.42) -- (.55,1.42) -- (pb.D);
 \draw (0,1.42) node[circ]{} -- (na.D); \draw (na.S) -- (nb.D);
 \draw (.55,1.42) -- ++(.55,0) node[right]{$Y$};
 \draw (nb.S) -- (0,-.99) node[below]{VSS};
 \draw (pa.G) -- ++(-.12,0) node[left]{$A$}; \draw (pb.G) -- ++(.12,0) node[right]{$B$};
 \draw (na.G) -- ++(-.12,0) node[left]{$A$}; \draw (nb.G) -- ++(-.12,0) node[left]{$B$};
 \node at (0,-1.64) {(b) NAND2};
\end{scope}
\begin{scope}[shift={(7.6,0)}]
 \node[pmos] (pb) at (0,2.35) {}; \node[pmos] (pa) at (0,1.32) {};
 \node[nmos] (na) at (-.55,-.05) {}; \node[nmos,xscale=-1] (nb) at (.55,-.05) {};
 \draw (pb.S) -- (0,3.02) node[above]{VDD}; \draw (pb.D) -- (pa.S);
 \draw (pa.D) -- (0,.64) node[circ]{}; \draw (na.D) -- (-.55,.64) -- (.55,.64) -- (nb.D);
 \draw (.55,.64) -- ++(.55,0) node[right]{$Y$};
 \draw (na.S) -- (-.55,-.81) -- (.55,-.81) -- (nb.S); \node[below] at (0,-.81) {VSS};
 \draw (pa.G) -- ++(-.12,0) node[left]{$A$}; \draw (pb.G) -- ++(-.12,0) node[left]{$B$};
 \draw (na.G) -- ++(-.12,0) node[left]{$A$}; \draw (nb.G) -- ++(.12,0) node[right]{$B$};
 \node at (0,-1.64) {(c) NOR2};
\end{scope}
\begin{scope}[shift={(11.6,0)}]
 \foreach \i/\gate in {0/A,1/B,2/C,3/D}{
   \node[pmos] (p\i) at ({1.02*\i},1.45) {};
   \draw (p\i.S) -- ({1.02*\i},2.02); \draw (p\i.D) -- ({1.02*\i},.85);
   \draw (p\i.G) -- ++(-.05,0) node[left,font=\scriptsize]{$\gate$};
 }
 \draw (0,2.02) -- (3.06,2.02); \node[above] at (1.53,2.02) {VDD};
 \draw (0,.85) -- (3.06,.85); \node[below] at (1.53,.85) {$Y$};
 \foreach \i/\gate in {0/A,1/B,2/C,3/D}{
   \node[nmos] (m\i) at (4.30,{2.1-.9*\i}) {};
   \draw (m\i.G) -- ++(-.05,0) node[left,font=\scriptsize]{$\gate$};
 }
 \draw (m0.D) -- ++(0,.12) node[above]{$Y$};
 \draw (m0.S) -- (m1.D); \draw (m1.S) -- (m2.D); \draw (m2.S) -- (m3.D);
 \draw (m3.S) -- ++(0,-.12) node[below]{VSS};
 \node at (1.95,-1.64) {(d) NAND4};
\end{scope}
\begin{scope}[shift={(19.0,0)}]
 \node[pmos,rotate=-90] (pg) at (0,1.8) {};
 \node[nmos,rotate=90] (ng) at (0,.5) {};
 \draw (pg.D) -- (-.85,1.8) -- (-.85,.5) -- (ng.D);
 \draw (pg.S) -- (.85,1.8) -- (.85,.5) -- (ng.S);
 \draw (-.85,1.13) node[circ]{} -- ++(-.3,0) node[left]{$X$};
 \draw (.85,1.13) node[circ]{} -- ++(.3,0) node[right]{$Z$};
 \draw (ng.G) -- ++(0,-.14) node[below]{$E$};
 \draw (pg.G) -- ++(0,.14) node[above]{$\overline E$};
 \node at (0,-1.64) {(e) TG};
\end{scope}
\end{circuitikz}}
\caption{Distinct gate topologies of the R1 AFC. For the NAND4 the two networks share output $Y$. The transmission gate (TG) connects $X$ and $Z$ when $E=1$. The latch adds the reset NMOS shown in Fig.~\ref{fig:r1afc-latch}.}
\label{fig:r1afc-cells}
\end{figure*}

The default inverter has $W_n/W_p=0.8/2.4~\mu\mathrm{m}$. Its latch and
clock-driver variants use $0.5/1.5$ and $8/24~\mu\mathrm{m}$.
NAND2, NOR2 and NAND4 use per-device widths of $1.6/2.4$,
$0.8/4.8$ and $3.2/2.4~\mu\mathrm{m}$, respectively. The TG uses
$1.2/3.6~\mu\mathrm{m}$. For an XOR, the generator forms
$U=\overline{AB}$ and then
$Y=\overline{\overline{AU}\,\overline{BU}}$. A two-input multiplexer
uses NAND2 cells and a complemented select. These same cells form the
Gray converters, incrementers, difference detectors and code update logic.

\begin{figure}[t]
\centering\resizebox{.98\columnwidth}{!}{%
\begin{circuitikz}[american,x=1cm,y=1cm,font=\PLLfigurefont\small,
 c/.style={draw,minimum width=1.12cm,minimum height=.66cm,align=center},
 ar/.style={->,>=latex}]
 \node[c] (di) at (1.05,1.05) {TG}; \node[c] (ii) at (3.65,1.05) {INV};
 \node[c] (io) at (5.5,1.05) {INV}; \node[c] (fb) at (3.65,2.7) {TG};
 \draw (-.37,1.05) node[left]{$D$} -- (di.west);
 \draw (di.east) -- (ii.west); \draw (ii.east) -- (io.west);
 \draw (io.east) -- (6.67,1.05) node[right]{$Q$};
 \draw (6.32,1.05) node[circ]{} |- (fb.east); \draw (fb.west) -| (2.32,1.05) node[circ]{};
 \draw (di.north) -- ++(0,.28) node[above]{open};
 \draw (fb.north) -- ++(0,.16) node[above]{hold};
 \draw (4.58,1.05) node[circ]{} -- ++(0,.32) node[above]{$\overline Q$};
 \node[nmos] (mr) at (2.32,-.12) {};
 \draw (mr.D) -- (2.32,1.05); \draw (mr.G) -- ++(-.14,0) node[left]{RST};
 \draw (mr.S) -- ++(0,-.12) node[below]{VSS};
 \node[anchor=west,font=\PLLfigurefont\scriptsize] at (2.45,.6) {$X$};
 \node[anchor=west,font=\PLLfigurefont\scriptsize] at (2.65,-.1) {$W_n=2~\mu$m};
 \begin{scope}[shift={(0,-2.55)}]
 \node[c] (ml) at (1.5,0) {Master latch};
 \node[c] (sl) at (4.8,0) {Slave latch};
 \draw[ar] (-.35,0) node[left]{$D$} -- (ml.west);
 \draw[ar] (ml.east) -- node[above]{$Q_M$} (sl.west);
 \draw[ar] (sl.east) -- (6.65,0) node[right]{$Q$};
 \node[below,align=center] at (ml.south) {open $L$\\hold $H$};
 \node[below,align=center] at (sl.south) {open $H$\\hold $L$};
 \end{scope}
\end{circuitikz}}
\caption{R1 resettable static latch (top) and ordinary positive-edge DFF (bottom). Each TG receives an enable and its complement. The feedback TG returns $Q$ to storage node $X$ during hold. Active reset pulls $X$ to VSS in each latch.}
\label{fig:r1afc-latch}
\end{figure}

Two latches form one DFF. The preset variant inverts $D$ and exposes
the slave's complementary output. Each shared clock group forms $K_d$
with six cascaded inverters. With active-high reset $R$, the interlock
computes
\begin{align}
 L &= \overline K\,\overline{K_d}\,\overline H\,\overline R, &
 H &= K\,K_d\,\overline L\,\overline R .
\end{align}
A NAND4 and two large inverters drive each phase and its complement.

R1 contains 146 DFFs and one additional latch on the feedback clock
gate. The 293 latches use 586 TGs, 586 latch inverters and 293 reset
NMOS devices. Outside these latches, the generator emits 531 inverters
(including clock drivers), 740 NAND2s, 79 NOR2s and 20 NAND4s. The
total is $293(9)+531(2)+740(4)+79(4)+20(8)=7{,}135$ MOS instances
in the AFC core. These counts can be checked against the selected R1
gate generator, its port manifest and the generated transistor netlist.


% IEEE appendix input. Compile inside the main document.
\section{Historical Circuit Variants}
\label{sec:historical}
Two early receiver families share a resistively loaded differential pair.
The load experiment exposes the two drain voltages directly. The
direct-coupled clock receiver adds two CMOS inverters at the positive
output (Fig.~\ref{fig:hist-rx}). In the selected receiver, a coupling
capacitor and a feedback resistor set the first inverter's input bias
(Fig.~\ref{fig:aprx}). This distinction matters when comparing the
measured loading and clock waveforms.

\begin{figure*}[!t]
\centering\resizebox{.90\textwidth}{!}{%
\begin{circuitikz}[american,font=\small]
% Direct CML: positive input drives the negative output drain.
\node[nmos] (hip) at (0,1.5) {}; \node[nmos,xscale=-1] (hin) at (2.8,1.5) {};
\node[nmos] (ht) at (1.4,-.15) {};
\draw (hip.S) -- (0,.66) -- (2.8,.66) -- (hin.S);
\draw (1.4,.66) node[circ]{} -- (ht.D);
\draw (ht.S) -- ++(0,-.18) node[below]{VSS};
\draw (ht.G) -- ++(-.15,0) node[left]{BN};
\draw (hip.G) -- ++(-.12,0) node[left]{INP};
\draw (hin.G) -- ++(.12,0) node[right]{INN};
\draw (hip.D) -- (0,2.25) node[circ]{} to[R] (0,3.9);
\draw (hin.D) -- (2.8,2.25) node[circ]{} to[R] (2.8,3.9);
\draw (0,3.9) -- (2.8,3.9) node[midway,above]{VDD};
\draw (0,2.25) -- ++(-.6,0) node[left]{ON};
\draw (2.8,2.25) -- ++(.55,0) node[right]{OP};
\node[align=center] at (1.4,-1.88) {(a) Shared CML front end};
% Repeated labels refer to the same OP drain. Do not connect OP through a capacitor.
\begin{scope}[xshift=6.4cm]
\node[american not port] (a) at (1.1,2.45) {};
\node[american not port] (b) at (3.9,2.45) {};
\draw (-.63,2.45) node[left]{OP} -- (a.in);
\draw (a.out) -- node[above]{$\overline{\mathrm{CLK}}$} (b.in);
\draw (b.out) -- ++(.3,0) node[right]{CLK};
\node[align=center] at (2.05,1.05) {Direct receiver extension\\OP drives the first inverter gate};
\node[align=center] at (2.00,-1.88) {(b) Two added CMOS inverters};
\end{scope}
\begin{scope}[xshift=13.5cm]
\node[nmos,rotate=90] (hs) at (1.3,2.75) {};
\draw (hs.D) -- ++(-.29,0) node[left]{CTRL};
\draw (hs.G) -- ++(0,-.13) node[below]{$s_{k+}$};
\draw (hs.S) -- (3.2,2.75) node[circ]{} -- ++(.43,0) node[right]{$g_k$};
\node[nmos] (hg) at (3.2,1.13) {};
\draw (hg.D) -- (3.2,2.75); \draw (hg.G) -- ++(-.1,0) node[left]{$\overline{s}_{k+}$};
\draw (hg.S) -- ++(0,-.19) node[below]{VSS};
\node[align=center] at (1.75,-1.88) {(c) Early fine-bank gate drive};
\end{scope}
\end{circuitikz}}
\caption{Archived load and oscillator variants. With $1$-k$\Omega$ loads the CML-only cell has four input fingers of either 1.5 or 3~$\mu$m per device. The direct receiver defaults to $3$-k$\Omega$ loads and adds (b). Both use a 12-finger tail with 2.5-$\mu$m fingers. In (c), each of $k=8,16,24$ uses an NMOS to pass CTRL to a supplemental sink gate and a second NMOS to discharge it. The later ring adds a parallel PMOS on the CTRL path.}
\label{fig:hist-rx}
\end{figure*}

\begin{figure*}[!t]
\centering\resizebox{.80\textwidth}{!}{%
\begin{circuitikz}[american,font=\small,
 tx/.style={draw,align=center,inner sep=3pt,minimum width=1.8cm,minimum height=.58cm}]
% Selector TG inputs are tied to named rail or mirror nodes.
\node[tx] (ub) at (1.7,3.35) {TG, UP};
\node[tx] (ur) at (1.7,2.40) {TG, $\overline{\rm UP}$};
\node[tx] (db) at (1.7,1.0) {TG, DN};
\node[tx] (dr) at (1.7,.05) {TG, $\overline{\rm DN}$};
\draw (ub.west) -- ++(-.42,0) node[left]{$b_{p,\rm cp}$};
\draw (ur.west) -- ++(-.42,0) node[left]{VDD};
\draw (db.west) -- ++(-.42,0) node[left]{$b_{n,\rm cp}$};
\draw (dr.west) -- ++(-.42,0) node[left]{VSS};
\draw (ub.east) -- (3.2,3.35) -- (3.2,2.40) -- (ur.east);
\draw (db.east) -- (3.2,1.0) -- (3.2,.05) -- (dr.east);
\node[pmos] (mu) at (5.15,2.86) {}; \node[nmos] (md) at (5.15,.53) {};
\draw (3.2,2.86) node[circ]{} -- node[above]{$g_p$} (mu.G);
\draw (3.2,.53) node[circ]{} -- node[above]{$g_n$} (md.G);
\draw (mu.S) -- ++(0,.32) node[above]{VDD}; \draw (md.S) -- ++(0,-.32) node[below]{VSS};
\draw (mu.D) -- (md.D); \draw (5.15,1.7) node[circ]{} -- ++(.5,0) node[right]{CTRL};
\begin{scope}[xshift=8.4cm]
\node[american not port] (iu) at (1.2,2.9) {};
\node[american not port] (idn) at (1.2,1.80) {};
\draw (iu.in) -- ++(-.3,0) node[left]{UP}; \draw (iu.out) -- ++(.3,0) node[right]{$\overline{\rm UP}$};
\draw (idn.in) -- ++(-.3,0) node[left]{DN}; \draw (idn.out) -- ++(.3,0) node[right]{$\overline{\rm DN}$};
\node[align=center] at (1.18,.30) {Mirror bias network\\as in Fig.~\ref{fig:apcp}};
\end{scope}
\end{circuitikz}}
\caption{Gate-switched charge-pump candidate. Each of the four TGs uses an NMOS/PMOS pair with complementary gate drives. The label gives the NMOS enable. With UP high the PMOS output gate selects the mirror voltage, and with UP low it selects VDD. DN likewise selects the NMOS mirror voltage or VSS. The TG and inverter widths are $0.5/1.5~\mu$m at $L=0.13~\mu$m. The NMOS and PMOS output widths are 4 and 12~$\mu$m at $L=1.04~\mu$m.}
\label{fig:hist-cp}
\end{figure*}

The original AFC core R0 contains 6,908 MOS instances and 145 DFFs.
Its FB counter uses the physical FB clock directly. Reset follows
$\mathrm{RST}\lor\overline{\mathrm{count\_run}}$. R1 adds a latched
FB-clock enable and holds counter reset deasserted during DRAIN,
as shown in Fig.~\ref{fig:r1afc-datapath}. The archive also records
the R2/counter9 proposal as a prepared design option. The nominal
closed-loop result in this report belongs to R1.

The simulator additionally accepts eight project-authored behavioral
modules. They represent a charge pump, precharge path, linear VCO,
divide-by-32 counter, PFD, search controller, handoff controller and
AFC unit-test source. Their equations and event handlers provide stimuli
and functional substitutes in staged benches. The transistor counts in
the selected design apply to the MOS implementations.

The companion \texttt{module-index.md} maps all 32 distinct Spectre
subcircuit names to these drawings and identifies the eight Verilog-A
modules by name.

\section{Conclusion}
The selected controller reached code 29 and acquired sustained nominal lock 1.119\,$\mu$s after fine enable. The weaker-clamp run exceeded its precharge tolerance even though the fine loop acquired. The fixed-code loop exhausted its analog tuning range within the observed interval.
\begin{thebibliography}{00}
\bibitem{student} \emph{A 3.2 GHz Fast-Locking and Low-Jitter PLL with Automatic Frequency Calibration Technique in 0.18-$\mu$m CMOS Technology}, supplied undergraduate design report, Figs.~1, 2, 4, 7, 9, 11 and 12.
\bibitem{jeong} C.-Y. Jeong, D.-H. Choi, and C. Yoo, \emph{A Fast Automatic Frequency Calibration (AFC) Scheme for Phase-Locked Loop (PLL) Frequency Synthesizer}, Proc. IEEE Radio Frequency Integrated Circuits Symp., 2009, pp.~583--586.
\bibitem{evidence} PLL130-AFC Project, \emph{Nominal Transient and Acquisition Evidence}, Oct. 2026. \url{https://github.com/handsomegary/pll130-afc}.
\end{thebibliography}
\end{document}
